\section{Explicit rate of convergence and rational seeds}
\label{sec:rate}

This section fixes the parameters of the construction of
Section~\ref{sec:extraction} as functions of the target index \(a\) and
adds up the bounds of Section~\ref{subsec:cost-bounds}, which gives the
rate of Theorem~\ref{thm:quantitative}. It then fixes every remaining
choice, identifies the field of coefficients
(Theorem~\ref{thm:effective}, Remark~\ref{rem:field-every-a}), and
passes to a rational seed (Corollary~\ref{cor:seed}); together these
results prove Theorems~\ref{thm:main} and~\ref{thm:main-relative}.

\subsection{An explicit rate of convergence}
\label{subsec:quantitative}

We bound the terms of the pooled exponent \eqref{eq:pooled-exponent} by
\eqref{eq:type-error} and \eqref{eq:finite-cost-bounds}. Without the
rounding loss \(\mathcal L\), the type penalty \(\mathcal T_N\) and the
three overheads, \(p\) would equal the profile exponent \(\tau_a\) of
\eqref{eq:tau}; we choose the parameters so that each of these
quantities is at most an absolute constant times \(K\). The rounding
loss is at most \(4J_a\) and the type penalty at \(N=1\) is
\(O(J_a\log K)\), with \(J_a=O(a^3)\), so the rounding scale
\(K=a^3\ell_a\) makes both \(O(K)\) and the type scale \(N=1\) suffices.
The copy overhead is \(O(I)=O(a^2)\). The interpolation overhead is at
most about twice the largest entropy charge of one relation, which is
\(O(K)\) already for \(k=1\); a larger interpolation power would lower
it but enlarge \(A=A_1^{kj}\). The catalyst has \(O(aK)\) factors, each
of rank at most \((5a)^3\), and the reuse count \(j=a\ell_a\) brings its
overhead down to \(O(K)\). Hence the numerator of \(p\) is
\(K(3\log a+O(1))\) and \(\Delta_1\ge K(\frac43\log a-O(1))\), so that
\(p-9/4=O(1/\log a)\), like \(\tau_a-9/4\). At small
\(a\), larger \(k\), \(j\) and \(K\) reduce the overheads further
(Section~\ref{sec:costs}, which keeps \(N=1\)).

\begin{theorem}[An explicit parameter family]\label{thm:quantitative}
Let \(a\ge2^{64}\) be an integer, and put
\[
 \ell_a=\lceil\log_2a\rceil,\qquad K=a^3\ell_a,\qquad N=k=1,
 \qquad j=a\ell_a.
\]
Take the occupations \(\lambda_{r,x},\mu_r\) of the row walk of
Proposition~\ref{prop:occupations}, in which the split at state \(x\)
has width \(h_{r,x}=\min(x,3r-x)\); round them at scale \(K\) as in
Section~\ref{sec:rounding}, by one of the roundings that
Lemma~\ref{lem:circulation-rounding} allows; and let
\(r_*\MM_A\restr m_*\MM_B\) be the relative scheme of
Proposition~\ref{prop:relative} built from these rounded counts. Then
\(B>A\), and the pooled exponent \(p\) and the source side \(A\) of this
scheme satisfy
\begin{equation}\label{eq:explicit-rate}
 \frac94<\tau_a\le p<\frac94+\frac{900}{\log a},
 \qquad
 \log A\le181a^5\ell_a^2.
\end{equation}
In particular, \(p\to9/4\) as \(a\to\infty\), and
\(\log A=O(a^5(\log a)^2)\).
\end{theorem}

\begin{proof}
We bound \(\Delta_1\) from below and each term of the numerator from
above, all divided by \(K\). Throughout, \(N=k=1\).

\emph{Elementary facts.} Since \(a\ge2^{64}\), we have
\(\log a\ge64\log2>32\) and \(\ell_a\ge64\); since
\(\log_2a\le\ell_a\le a\), also \(\log K\le4\log a\) and
\((\log a)/\ell_a\le\log2\). We use the numerical bounds
\(\log2<7/10\), \(\log4<2\), \(\log5<2\) and \(\log181<6\). The number
of transition coordinates satisfies \(J_a\le2a^3\) by
\eqref{eq:coordinate-count}. The number of relations satisfies
\(I\le3(a+2)(a-1)/2\le2a^2\), because row \(r\) has \(3r-1\) interior
states and one tripling. Finally, by \eqref{eq:rounded-edge-errors},
\(w_{r,x}<K\lambda_{r,x}+d_{r,x}\) and \(z_r<K\mu_r+1\), where
\(\lambda_{r,x}<180\) and \(\mu_r\le2\) by
Proposition~\ref{prop:occupations}, and
\(d_{r,x}\le\lfloor3a/2\rfloor+1\le2a\). Hence \(w_i<180K+2a\) for
every split \(i\) and \(z_r<2K+1\) for every tripling \(r\).

\emph{Step 1: a lower bound for \(\Delta_1\).} We show that
\(\Delta_1/K>\frac43\log a-7>0\). The type penalty
\eqref{eq:type-penalty} at \(N=1\) satisfies
\begin{equation}\label{eq:quantitative-type-loss}
 \mathcal T_1\le J_a\log(180K+2a+1)+(a-1)\log(2K+2).
\end{equation}
Indeed, the coefficient of \(\log(w_i+1)\) in \(\mathcal T_1\) is the
number of children of split \(i\) with positive count, minus one; it is
less than \(d_i\), so these coefficients sum to at most \(J_a\). There
are at most \(a-1\) triplings, and each contributes
\(\log(z_r+1)\le\log(2K+2)\). Since \(180K+2a+1\le181K\) and
\(2K+2\le4K\), the two logarithms in \eqref{eq:quantitative-type-loss}
are less than \(6+4\log a\) and \(2+4\log a\), respectively. Combining
this with the bound \(\mathcal L\le4J_a\) on the rounding loss
(Proposition~\ref{prop:entropy-loss}), with \(J_a\le2a^3\), and with
\(K=a^3\ell_a\), we obtain
\[
 \frac{4J_a+\mathcal T_1}{K}
 \le\frac{20+8\log a+\frac{a-1}{a^3}(2+4\log a)}{\ell_a}
 \le\frac{\frac{41}2+9\log a}{\ell_a}
 <\frac{41}{128}+\frac{63}{10}.
\]
The middle inequality uses \((a-1)/a^3\le1/a^2\le1/4\); the last uses
\(\ell_a\ge64\) and \((\log a)/\ell_a\le\log2<7/10\).

The inequality
\(\log(1+\xi)\ge\xi-\xi^2/2\) for \(\xi\ge0\), applied with
\(\xi=1/(3h)\), together with \(\sum_{h=1}^{a-1}1/h\ge\log a\) and
\(\sum_{h\ge1}1/h^2<2\), gives
\[
 \log s_a
 \ge\frac13\sum_{h=1}^{a-1}\frac1h
     -\frac1{18}\sum_{h=1}^{a-1}\frac1{h^2}
 \ge\frac13\log a-\frac19.
\]
Since \(F(a,a)=s_a(3a-1)/2\ge as_a\), it follows that
\(\log F(a,a)\ge\frac43\log a-\frac19\). By \eqref{eq:type-error} with
\(N=1\) and by Proposition~\ref{prop:entropy-loss},
\(\Delta_1\ge K\log F(a,a)-\mathcal L-\mathcal T_1\). Therefore
\begin{equation}\label{eq:quantitative-contraction}
 \frac{\Delta_1}{K}
 \ge\log F(a,a)-\frac{4J_a+\mathcal T_1}{K}
 >\frac43\log a-7>0,
\end{equation}
because \(\frac19+\frac{41}{128}+\frac{63}{10}<7\) and \(\log a>32\).
As \(\log(B/A)=kj\Delta_1\), this proves \(B>A\).

\emph{Step 2: copy overhead.} By \eqref{eq:finite-cost-bounds} and
\(I\le2a^2\),
\[
 \frac{\log(L/M^3)}{K}<\frac{3I\log2}{K}\le\frac{6}{a\ell_a}\le1.
\]

\emph{Step 3: interpolation overhead.} By \eqref{eq:finite-cost-bounds},
\(\log E\le\log(3I+1)+2\bar H_1\), where \(\bar H_1=\max_iH_i\) is the
largest entropy charge. For a split, \(H_i\le w_i\log d_i\), because its
empirical child distribution has at most \(d_i\) outcomes, and
\(w_i\log d_i\le180K+2a\log(2a)\) by \eqref{eq:rounded-count-bounds}.
For a tripling, \(H_r=z_r\log3<(2K+1)\log3\), which is smaller. Hence
\(\bar H_1\le180K+2a\log(2a)\), and
\[
 \frac{\log E}{K}
 \le360+\frac{\log(6a^2+1)+4a\log(2a)}{K}
 \le360+\frac{2a+4a^2}{a^3\ell_a}
 \le360+\frac{5}{a\ell_a}\le361,
\]
where we used \(\log(6a^2+1)\le2a\) and \(\log(2a)\le a\), both valid
for \(a\ge2\).

\emph{Step 4: catalyst overhead.} By \eqref{eq:catalyst-size}, the
catalyst has \(\bar w+\bar z-K\) factors. The count bounds
\eqref{eq:rounded-count-bounds}, together with \(\Gamma_a\le180a\) and
\(J_a\le2a^3\), give
\[
 \bar w+\bar z-K\le180aK+2a^3+(2K+1)(a-1)-K
 \le182aK+2a^3+a<184aK.
\]
Hence \eqref{eq:finite-cost-bounds} gives
\(\log\kappa_1\le3\cdot184aK\log(5a)\), and with \(j=a\ell_a\),
\[
 \frac{\log\kappa_1}{jK}
 \le\frac{552\log(5a)}{\ell_a}
 \le552\log2\left(1+\frac{\log5}{\log a}\right)
 <552\cdot\frac7{10}\left(1+\frac28\right)=483.
\]
The middle inequality uses \(\ell_a\ge(\log a)/\log2\); the last uses
\(\log2<7/10\), \(\log5<2\) and \(\log a>8\).

\emph{Step 5: combination.} Since \(3\log(2a-1)<3\log a+3\), Steps~2--4
bound the numerator of \(p\) by
\(K(3\log a+3+1+361+483)<K(3\log a+850)\). Together with Step~1, this
gives
\[
 p<\frac{3\log a+850}{\frac43\log a-7}
 =\frac94+\frac{850+\frac{63}4}{\frac43\log a-7}
 <\frac94+\frac{900}{\log a},
\]
because \(\frac43\log a-7\ge\log a\) when \(\log a\ge21\). The
fixed-batch bound \(\Phi\) of \eqref{eq:finite-exponent} obeys the same
bound, since the extra term \((\log2)/(kj)<1\le K\) of its numerator
still leaves it below \(K(3\log a+850)\). The lower bound
\(p\ge\tau_a>9/4\) holds for every scheme of
Section~\ref{sec:extraction} with \(B>A\) (see \eqref{eq:tau}), hence
here by Step~1.

\emph{Step 6: source side.} With \(N=1\), \(\log A_1=W\) is the sum of
the logarithms of the auxiliary dimensions, weighted by the split
counts. By \eqref{eq:rounded-count-bounds},
\(W\le K\Gamma_a+J_a\log(2a)\le180a^4\ell_a+2a^3\log(2a)\), so
\[
 \log A=kj\log A_1
 \le a\ell_a\bigl(180a^4\ell_a+2a^3\log(2a)\bigr)
 =a^5\ell_a^2\left(180+\frac{2\log(2a)}{a\ell_a}\right)
 \le181a^5\ell_a^2,
\]
since \(\log(2a)\le a\) and \(\ell_a\ge64\). The limit statements follow
because \(\ell_a=O(\log a)\).
\end{proof}

\begin{corollary}[An index for a given allowance]
\label{cor:explicit-allowance}
Let \(\eps\) be a positive rational, and put
\[
 \ell=\max\left\{64,\left\lceil\frac{1800}{\eps}\right\rceil\right\},
 \qquad a=2^\ell.
\]
The relative scheme of Theorem~\ref{thm:quantitative} for this \(a\)
has pooled exponent \(p<9/4+\eps\), and its source side satisfies
\(\log A\le181\cdot2^{5\ell}\ell^2\). The index \(a\) and the
parameters \(K,N,k,j\) are computed from \(\eps\) by integer and
rational arithmetic.
\end{corollary}

\begin{proof}
Here \(a\ge2^{64}\) and \(\ell_a=\ell\). Since
\(\log a=\ell\log2>\ell/2\), we get
\(900/\log a<1800/\ell\le\eps\), and \eqref{eq:explicit-rate} gives
both claims.
\end{proof}

\begin{remark}[Calibration]\label{rem:calibration}
The family of Theorem~\ref{thm:quantitative} is chosen to keep the
proof short, and its bound is weak for moderate \(a\): at \(a=2^{64}\),
\eqref{eq:explicit-rate} gives only \(p<22.54\); the bound
\(9/4+900/\log a\) is below \(3\) exactly when \(\log a>1200\); and for
\(\eps=1/10\), Corollary~\ref{cor:explicit-allowance} takes
\(a=2^{18000}\).
Section~\ref{sec:costs} obtains much better bounds for small \(a\)
(Table~\ref{tab:finite-multiway}).
\end{remark}

\subsection{An effective algorithm over a cyclotomic field}
\label{subsec:effective}

Theorem~\ref{thm:quantitative} allows every rounding of
Lemma~\ref{lem:circulation-rounding}. To obtain one algorithm for each
allowance, we fix every choice by a rule and record which roots of unity
the maps use.

\begin{theorem}[Effective construction]\label{thm:effective}
There is a terminating procedure that, given a positive rational
\(\eps\), specifies a finite arithmetic algorithm that multiplies
arbitrary complex \(n\times n\) matrices in \(O_\eps(n^{9/4+\eps})\)
operations. All coefficients of the algorithm lie in a cyclotomic field
\(\Q(\zeta_\Lambda)\) whose order \(\Lambda\) the procedure computes.
The procedure uses only finite integer and rational computations,
coordinate permutations, tensor products, and the maps of
Sections~\ref{sec:preliminaries} and~\ref{sec:maps}--\ref{sec:extraction}.
\end{theorem}

\begin{proof}
We list the rules in the order in which they are applied, and then
collect the roots of unity.

\emph{Step 1: parameters.} Compute \(\ell\) and \(a\) from \(\eps\) by
Corollary~\ref{cor:explicit-allowance}, and \(K,N,k,j\) by
Theorem~\ref{thm:quantitative}.

\emph{Step 2: occupations and rounded counts.} Solve the rational linear
systems of Proposition~\ref{prop:occupations} by exact Gaussian
elimination. Round each row separately by
Lemma~\ref{lem:circulation-rounding}, using breadth-first (shortest)
augmenting paths. List the arcs in the following order, which is also
the order in which the search scans the arcs at each vertex (the reverse
residual arc of an arc is scanned at that arc's position): the
transition edges by increasing state \(x\) and then by increasing child
\(y\); then the tripling edge \(3r\to1\); then the source and sink arcs
by increasing state. From the rounded counts \(n_{r,x,y}\) and \(z_r\),
form the ledger and the catalyst multiplicities \(c_{\sigma,r,b}\) of
Proposition~\ref{prop:rounded-ledger}. The examples in
Section~\ref{sec:costs} use the same rule.

\emph{Step 3: maps.} Order tensor positions and sector words
lexicographically. Use the determinant bases
\eqref{eq:determinant-basis}, the interval sectors of
Lemma~\ref{lem:tripling}, the Fourier substitutions of
Proposition~\ref{prop:separator}, the rank decompositions
\eqref{eq:convolution-rank}, and the root-of-unity averages of
Lemma~\ref{lem:exactification}. Separate the sectors before applying
their individual coordinate maps, and compose the maps in the order of
the proofs of Propositions~\ref{prop:catalytic}
and~\ref{prop:relative}; Appendix~\ref{sec:composed-maps} writes out
the composed relation maps. This gives the relative scheme of
Proposition~\ref{prop:relative}.

\emph{Step 4: the field of the scheme.} The maps use roots of unity of
the following orders only:
\(E=kD+1\), in the exactification of Lemma~\ref{lem:exactification};
\(2a-1\), in the rank decomposition \eqref{eq:convolution-rank} of
\(C(a,a)\), through which Proposition~\ref{prop:relative} replaces each
source copy of \(S^+_{a,a}\) by \((2a-1)^3\) rank-one tensors;
\(2M_i-1\) for \(i\in\mathcal J\), in the separator of relation \(i\);
and \(r+b-1\) for every catalyst factor \(S^\sigma_{r,b}\) with
\(c_{\sigma,r,b}>0\), in the rank decompositions that supply the
catalyst \eqref{eq:kappa}. Every other map has rational coefficients:
the determinant bases, the interval sectors, the coordinate
permutations, the scalars \(1/E\) and \(1/(2M_i-1)\), and the
restriction of the output catalyst to a coefficient equal to one;
catalyst reuse introduces no further roots of unity. Hence all
coefficients of the scheme lie in \(\Q(\zeta_\Lambda)\), where
\begin{equation}\label{eq:cyclotomic-order}
 \Lambda=\lcm\bigl(E,\;2a-1,\;\{2M_i-1:i\in\mathcal J\},\;
       \{r+b-1:c_{\sigma,r,b}>0\}\bigr).
\end{equation}
We represent \(\Q(\zeta_\Lambda)\) by the power basis modulo the
\(\Lambda\)-th cyclotomic polynomial.

\emph{Step 5: the algorithm.} By Theorem~\ref{thm:quantitative} and
Corollary~\ref{cor:explicit-allowance}, \(2<p<9/4+\eps\), so
Lemma~\ref{lem:batch} turns the scheme into an algorithm with
\(O(n^p)\subseteq O_\eps(n^{9/4+\eps})\) operations. Its zero padding
and its recursion introduce no further roots of unity, so all its
coefficients lie in \(\Q(\zeta_\Lambda)\). The procedure thus specifies
finite coefficient arrays and the order of their composition.
\end{proof}

\begin{remark}[The field for every index]\label{rem:field-every-a}
Steps~2--4 of this proof do not use the special form \(a=2^\ell\) of
Step~1. Hence, for every integer \(a\ge2^{64}\), the relative scheme of
Theorem~\ref{thm:quantitative}, built with the rounding rule of Step~2
and the orders of Step~3, is produced from \(a\) by a terminating
procedure, and its coefficients lie in \(\Q(\zeta_\Lambda)\), with
\(\Lambda\) given by \eqref{eq:cyclotomic-order}. Together with
Theorem~\ref{thm:quantitative}, this proves
Theorem~\ref{thm:main-relative}.
\end{remark}

\subsection{Rational coefficients and a square seed}
\label{subsec:seed}

\begin{lemma}[Restriction of coefficients to the rationals]
\label{lem:descent}
Let \(T\) be a tensor with rational coefficients, and suppose that a
decomposition of \(T\) into \(R\) rank-one tensors is given over a
number field \(\mathbb K\) of degree \(g=[\mathbb K:\Q]\). Then the
construction below gives a decomposition of \(T\) into at most \(g^2R\)
rank-one tensors with rational coefficients.
\end{lemma}

\begin{proof}
Choose a basis \(\theta_1=1,\theta_2,\ldots,\theta_g\) of \(\mathbb K\)
over \(\Q\) (for \(\Q(\zeta_\Lambda)\), the power basis, in which all
field operations are effective), and let \(\varphi\colon\mathbb K\to\Q\)
be the \(\Q\)-linear map that extracts the coefficient of \(\theta_1\).
In one summand of the decomposition, write the first two linear forms as
\(\sum_s\theta_s\alpha_s\) and \(\sum_{s'}\theta_{s'}\beta_{s'}\), where
the forms \(\alpha_s\) and \(\beta_{s'}\) have rational coefficients,
and call the third form \(\gamma\). Apply \(\varphi\) to every
coefficient of the tensor identity. The summand becomes
\[
 \sum_{s,s'=1}^{g}\alpha_s\otimes\beta_{s'}\otimes
 \varphi(\theta_s\theta_{s'}\gamma),
\]
where \(\varphi\) acts on the coefficients of \(\gamma\). This is a sum
of at most \(g^2\) rank-one tensors with rational coefficients; terms
with \(\alpha_s=0\) or \(\beta_{s'}=0\) are omitted, so a summand with
rational coefficients is left unchanged. The tensor \(T\) itself is
unchanged, because its coefficients are rational and \(\varphi\) is the
identity on \(\Q\).
\end{proof}

\begin{corollary}[Rational seeds]\label{cor:seed}
For every positive rational \(\eps\), a terminating procedure produces
an integer \(d_0\ge2\) and a decomposition of \(\MM_{d_0}\) into
\(R_{d_0}<d_0^{9/4+\eps}\) rank-one tensors with rational coefficients.
Consequently, a fixed rational bilinear recursion attains an exponent
strictly below \(9/4+\eps\).
\end{corollary}

\begin{proof}
Apply Theorem~\ref{thm:effective} with allowance \(\eps/2\). For each
side \(n\), unroll the pooled recursion of Lemma~\ref{lem:batch} down to
its classical leaves, with the zero padding of that lemma. By the last
assertion of Lemma~\ref{lem:batch}, every leaf multiplication
multiplies a linear form in the entries of the left input by a linear
form in the entries of the right input. The unrolled algorithm is
therefore a decomposition of \(\MM_n\) into rank-one tensors, one for
each leaf multiplication. Their number is at most the arithmetic cost,
hence \(O_\eps(n^{9/4+\eps/2})\), and their coefficients lie in
\(\Q(\zeta_\Lambda)\), a field that does not depend on \(n\). Apply
Lemma~\ref{lem:descent}, and let \(R_n\) be the number of rank-one
tensors after descent. Since the degree
\(g=[\Q(\zeta_\Lambda):\Q]\) does not depend on \(n\), we have
\(R_n\le g^2\cdot O_\eps(n^{9/4+\eps/2})=O_\eps(n^{9/4+\eps/2})\).

Write \(9/4+\eps=\alpha_\eps/\beta_\eps\) with positive integers
\(\alpha_\eps,\beta_\eps\). For \(n=2,3,\ldots\), compute \(R_n\) and
stop at the first \(n\) with \(R_n^{\beta_\eps}<n^{\alpha_\eps}\); this
\(n\) is \(d_0\). The loop terminates because the exponent of the bound
on \(R_n\) is smaller than \(9/4+\eps\) by \(\eps/2\). Ordinary block
recursion on the seed has exponent \(\log_{d_0}R_{d_0}<9/4+\eps\); its
fixed linear maps cost \(O(n^2)\) operations per node.
\end{proof}

\begin{remark}[Size of the seed]\label{rem:seed-size}
Suppose that \(\eps\le3/4\), the only case of interest, since classical
multiplication has exponent \(3\). Let \(A\) and \(B\) be the sides of
the scheme that the proof of Corollary~\ref{cor:seed} uses, for the
allowance \(\eps/2\). Every side \(n\le A\) lies below the leaf
threshold \(\lceil A/(1-A/B)\rceil+1\) of Lemma~\ref{lem:batch}, so for
such \(n\) the unrolled algorithm is the classical one, with
\(n^3\ge n^{9/4+\eps}\) rational products, which descent leaves
unchanged; hence \(d_0>A\). Moreover \(A=A_1^{kj}=A_1^j\ge2^j\) with
\(j=a\ell_a\ge2^{70}\), so \(d_0>2^{2^{70}}\); here \(A_1\ge2\) because,
by \eqref{eq:rounded-row-flow}, the splits that enter state \(0\) of a
row have total count \(K\), so some split is active. Since the seed has
rational coefficients, its identity holds over every field of
characteristic zero. Theorem~\ref{thm:bounded-seed} gives an explicit
upper bound for the side of a seed obtained from the rational variant of
the construction.
\end{remark}
